\section*{Independent reconstruction of B06-A028-PAID-U-BOUNDARY-v1}

\paragraph{Audited claim.}
Let $f$ be Boolean and let $\mathcal T$ be a finite physical restriction
tree.  At an internal node $v$, reached with probability $p_v$, the live
coordinate $i_v$ is queried and $f_v$ is the restricted Boolean function.
Write
\[
 \rho_i(g)=\Delta_iU(g)-K_i^{(0)}(g)
           =B_i(g)-\Gamma_i(g)+K_i^{(+)}(g),
\]
\[
 \Pi_I(\mathcal T)=\sum_{v\in\mathcal T^\circ}
                 p_v\operatorname{Inf}_{i_v}(f_v).
\]
The source theorem assumes fixed $A,\kappa\ge0$ and $0\le b<1$, a tree
for every nonconstant $f$, leaf estimates
\[
 U(f_\ell)\le\kappa I(f_\ell),                              \tag{1}
\]
and the aggregate residual estimate
\[
 \sum_vp_v\rho_{i_v}(f_v)\le A\Pi_I(\mathcal T)+bU(f).       \tag{2}
\]
It concludes
\[
 U(f)\le{\max\{A,\kappa\}I(f)+(2/\ln2)\operatorname{Var}(f)\over1-b}.
                                                                    \tag{3}
\]

\paragraph{Exact telescope audit.}
At a single live query,
\[
 U(g)=\Delta_iU(g)+\frac{U(g_{i=1})+U(g_{i=-1})}{2}.         \tag{4}
\]
Multiplying by the reach probability and summing cancels every internal
child exactly once, giving
\[
 \sum_vp_v\Delta_{i_v}U(f_v)
 =U(f)-\sum_\ell p_\ell U(f_\ell).                          \tag{5}
\]
Fourier restriction gives
\[
 I(g)=\operatorname{Inf}_i(g)
      +\frac{I(g_{i=1})+I(g_{i=-1})}{2},                    \tag{6}
\]
so the same cancellation yields
\[
 \Pi_I(\mathcal T)
 =I(f)-\sum_\ell p_\ell I(f_\ell).                          \tag{7}
\]
This is an equality, not a comparison: queried influence and terminal
influence are disjoint in the root ledger.

\paragraph{Independent root-row reconstruction.}
At node $v$, let
\[
 a_v=\mathbb E[f\mid E_v],\qquad
 d_v=\mathbb E[f x_{i_v}\mid E_v],\qquad m_v=a_v^2+d_v^2.
\]
Boolean convolution gives the empty-row formula
\[
 K_{i_v}^{(0)}(f_v)
 =m_v\left[1-h_2\left({1+2a_vd_v/m_v\over2}\right)\right].
\]
The scalar inequality
\[
 1-h_2((1+x)/2)\le{1-\sqrt{1-x^2}\over\ln2}
\]
therefore gives
\[
 K_{i_v}^{(0)}(f_v)
 \le {m_v-|a_v^2-d_v^2|\over\ln2}
 ={2\min(a_v^2,d_v^2)\over\ln2}
 \le {2d_v^2\over\ln2}.                                    \tag{8}
\]
The node functions
\[
 \psi_v={\mathbf1_{E_v}x_{i_v}\over\sqrt{p_v}}
\]
are orthonormal.  Incomparable nodes have disjoint supports; at a strict
descendant, the ancestor coordinate is fixed and the new query coordinate
is still conditionally fair.  Since
$\langle f,\psi_v\rangle=\sqrt{p_v}d_v$, Bessel gives
\[
 \sum_vp_vK_{i_v}^{(0)}(f_v)
 \le {2\over\ln2}\sum_vp_vd_v^2
 \le {2\over\ln2}\operatorname{Var}(f).                    \tag{9}
\]
The same physical coordinate may appear at incomparable nodes without
affecting orthogonality.  It cannot reappear on one path because a physical
restriction has fixed it.

\paragraph{Assembly.}
Using $\Delta_iU=K_i^{(0)}+\rho_i$, equations (2), (5), and (9) give
\[
 U(f)-\sum_\ell p_\ell U(f_\ell)
 \le {2\over\ln2}\operatorname{Var}(f)+A\Pi_I+bU(f).
\]
Apply (1) and (7):
\[
 (1-b)U(f)
 \le {2\over\ln2}\operatorname{Var}(f)+AI(f)
       +(\kappa-A)\sum_\ell p_\ell I(f_\ell).
\]
Because $0\le\sum_\ell p_\ell I(f_\ell)\le I(f)$, this is (3).
The source reduction $H\le2I+2U$ then gives
\[
 H(f)\le
 \left(2+{2(\max\{A,\kappa\}+2/\ln2)\over1-b}\right)I(f).   \tag{10}
\]
Constants have $I=U=0$ and are separate endpoints.

\paragraph{Verdict and first unsupported universal line.}
The proof HOLDS at its stated conditional scope.  All quantifiers,
normalizations, signs, terminal costs, and the single-use influence ledger
are valid.  It is not yet a universal GU proof.  The first unpaid line is
the existence, for every Boolean $f$, of one finite physical tree satisfying
(1)--(2) with a single fixed triple $(A,\kappa,b)$, $b<1$.  The previous
LTF-terminal modulation theorem supplies (1) only for a proper family and
does not prove (2) for an arbitrary non-flat shared-data collision.

The source theorem is an explicitly sufficient global $U$-drift route to
FEI, not a proof of $G\le aI+bU$ itself.  Relabeling (2) as universal GU
would be a quantifier error.
